\documentclass[tikz,border=4pt]{standalone}
\usepackage{ctex}
\usepackage{amsmath}
\usetikzlibrary{calc, arrows.meta, decorations.markings}

\begin{document}
% thinking-toolkit/06 算两次思想：直角三角形面积两种表达
% 直角三角形 ABC，∠C=90°，AC=6，BC=8，AB=10
% 面积法1：(1/2)*AC*BC；面积法2：(1/2)*AB*h
\begin{tikzpicture}[scale=0.9, thick]
  % 顶点坐标
  \coordinate (A) at (0,6);
  \coordinate (B) at (8,0);
  \coordinate (C) at (0,0);

  % 直角标记
  \draw (C) ++(0.3,0) -- ++(0,0.3) -- ++(-0.3,0);

  % 三角形
  \draw[thick] (A) -- (B) -- (C) -- cycle;

  % 顶点标签
  \node[left, font=\small] at (A) {$A$};
  \node[right, font=\small] at (B) {$B$};
  \node[below left, font=\small] at (C) {$C$};

  % 边长标注
  \node[left, font=\small] at ($(A)!0.5!(C)$) {$AC=6$};
  \node[below, font=\small] at ($(B)!0.5!(C)$) {$BC=8$};
  \node[above right, font=\small] at ($(A)!0.5!(B)$) {$AB=10$};

  % C到AB的高（垂足H）
  % H = 投影：H = C + t*(B-A)，t = dot(C-A, B-A)/|B-A|^2
  % B-A = (8,-6), C-A = (0,-6)
  % t = (0*8+(-6)*(-6))/(64+36) = 36/100 = 0.36
  % H = (0+0.36*8, 6+0.36*(-6)) = (2.88, 3.84)
  \coordinate (H) at (2.88, 3.84);
  \draw[red, dashed] (C) -- (H);
  \fill[red] (H) circle(2pt);
  \node[right, red, font=\small] at (H) {$H$};
  % 垂足标记
  \begin{scope}
    \clip ($(H)+(-0.5,-0.5)$) rectangle ($(H)+(0.5,0.5)$);
    \draw[red] ($(H)!0.25!(C)$) -- ($(H)!0.25!(C)+(0.22,0.17)$) -- ($(H)+(0.22,0.17)$);
  \end{scope}

  % 高的标注
  \node[below left, red, font=\small] at ($(C)!0.5!(H)$) {$h$};

  % 两种面积表达式
  \node[font=\small, align=center] at (4,-1.0) {
    方法1：$S = \dfrac{1}{2}\times AC \times BC = \dfrac{1}{2}\times 6\times 8 = 24$
  };
  \node[font=\small, align=center] at (4,-1.7) {
    方法2：$S = \dfrac{1}{2}\times AB \times h = \dfrac{1}{2}\times 10\times h$
  };
  \node[red, font=\small, align=center] at (4,-2.35) {
    两式相等 $\Rightarrow\ h = \mathbf{4.8}$
  };
\end{tikzpicture}
\end{document}
